\documentclass[11pt,a4paper]{article}
\usepackage[utf8]{inputenc}
\usepackage[T1]{fontenc}
\usepackage{lmodern}
\usepackage[french]{babel}
\usepackage{amsmath,amssymb,mathtools}
\usepackage{icomma}
\usepackage{xcolor}
\usepackage{colortbl}
\usepackage{tikz}
\usepackage[most]{tcolorbox}
\usepackage{enumitem}
\usepackage{array}
\usepackage{booktabs}
\usepackage{fancyhdr}
\usepackage[hidelinks]{hyperref}
\usetikzlibrary{arrows.meta,positioning,calc,patterns,backgrounds,shapes.geometric,decorations.pathreplacing,decorations.markings}
\usepackage[a4paper,margin=2.15cm,headheight=26pt,headsep=12pt]{geometry}
\usepackage{titlesec}

% ---------------- Palette ----------------
\definecolor{primary}{RGB}{0,151,169}
\definecolor{primarydark}{RGB}{0,88,101}
\definecolor{defcol}{RGB}{0,114,178}
\definecolor{thmcol}{RGB}{0,140,120}
\definecolor{formulacol}{RGB}{198,93,9}
\definecolor{excol}{RGB}{0,135,90}
\definecolor{attncol}{RGB}{200,40,55}
\definecolor{methcol}{RGB}{90,90,100}
\definecolor{remarkcol}{RGB}{120,94,200}
\definecolor{barcol}{RGB}{0,151,169}
\definecolor{barcold}{RGB}{0,110,124}
\definecolor{modecol}{RGB}{200,55,55}
\definecolor{meancol}{RGB}{0,90,200}
\definecolor{medcol}{RGB}{160,90,190}
\definecolor{quartcol}{RGB}{0,135,90}
\definecolor{ogivecol}{RGB}{176,74,0}

% ---------------- Boîtes ----------------
\tcbset{boxbase/.style={enhanced,breakable,boxrule=0.9pt,arc=2.5pt,
  left=3mm,right=3mm,top=2.3mm,bottom=2.3mm,
  fonttitle=\bfseries,coltitle=white,toptitle=0.6mm,bottomtitle=0.6mm}}

\newtcolorbox{definition}[1][]{boxbase,colback=defcol!6,colframe=defcol,
  title={\textsc{Définition}\ifx\relax#1\relax\relax\else\textmd{ : \itshape #1}\fi}}
\newtcolorbox{propriete}[1][]{boxbase,colback=thmcol!7,colframe=thmcol,
  title={\textsc{Propriété}\ifx\relax#1\relax\relax\else\textmd{ : \itshape #1}\fi}}
\newtcolorbox{formule}[1][]{boxbase,colback=formulacol!7,colframe=formulacol,
  title={\textsc{Formule}\ifx\relax#1\relax\relax\else\textmd{ : \itshape #1}\fi}}
\newtcolorbox{exemple}[1][]{boxbase,colback=excol!6,colframe=excol,
  title={\textsc{Exemple}\ifx\relax#1\relax\relax\else\textmd{ : \itshape #1}\fi}}
\newtcolorbox{methode}[1][]{boxbase,colback=methcol!8,colframe=methcol,sharp corners,
  title={\textsc{Méthode}\ifx\relax#1\relax\relax\else\textmd{ : \itshape #1}\fi}}
\newtcolorbox{remarque}[1][]{boxbase,colback=remarkcol!7,colframe=remarkcol,
  title={\textsc{Astuce}\ifx\relax#1\relax\relax\else\textmd{ : \itshape #1}\fi}}
\newtcolorbox{attention}[1][]{boxbase,colback=attncol!7,colframe=attncol,
  title={\textsc{Attention}\ifx\relax#1\relax\relax\else\textmd{ : \itshape #1}\fi}}

% Boîte "exercice" et "corrigé" (titre libre passé en argument)
\newtcolorbox{exobox}[1]{boxbase,colback=primary!4,colframe=primary,
  before skip=8pt,after skip=8pt,title={#1}}
\newtcolorbox{corbox}[1]{boxbase,colback=excol!5,colframe=excol!85!black,
  before skip=8pt,after skip=8pt,title={#1}}

% ---------------- En-tête ----------------
\newcommand{\docpart}[1]{\def\thedocpart{#1}}
\docpart{}
\pagestyle{fancy}
\fancyhf{}
\fancyhead[L]{\small\textcolor{primarydark}{\textbf{Fiche 8} — Statistiques descriptives}}
\fancyhead[R]{\small\textcolor{primarydark}{\textsc{\thedocpart}}}
\fancyfoot[C]{\small\thepage}
\renewcommand{\headrulewidth}{0.8pt}
\renewcommand{\headrule}{\hbox to\headwidth{\color{primary}\leaders\hrule height \headrulewidth\hfill}}

\titleformat{\section}{\large\bfseries\color{primarydark}}{\thesection}{0.6em}{}
\titleformat{\subsection}{\normalsize\bfseries\color{primary}}{\thesubsection}{0.6em}{}
\setlist{topsep=2pt,itemsep=1.5pt,parsep=0pt}

% Bandeau de titre du document
\newcommand{\bandeau}[2]{%
\begin{tcolorbox}[boxbase,colback=primary,colframe=primarydark,halign=center,
   before skip=2pt,after skip=12pt]
{\LARGE\bfseries\color{white} #1}\\[3pt]{\normalsize\color{white} #2}
\end{tcolorbox}}

\newcommand{\ecm}{\sigma_m}
\newcommand{\ect}{\sigma}
\newcommand{\med}{M\!e}
\docpart{Moyen — Corrigé — Thème 4}
\begin{document}
\bandeau{Thème 4 — Les indicateurs de dispersion}{Niveau Moyen — corrigé}

\begin{corbox}{Corrigé 1 — Écart moyen}
\begin{center}
\renewcommand{\arraystretch}{1.15}
\rowcolors{2}{primary!5}{white}
\begin{tabular}{|c|c|c|c|}
\hline
\rowcolor{primary!18}
$x_i$ & $n_i$ & $|m-x_i|$ & $n_i|m-x_i|$ \\\hline
2  & 4  & 8{,}8 & 35{,}2 \\
6  & 8  & 4{,}8 & 38{,}4 \\
10 & 10 & 0{,}8 & 8{,}0 \\
14 & 12 & 3{,}2 & 38{,}4 \\
18 & 6  & 7{,}2 & 43{,}2 \\\hline
\multicolumn{3}{|r|}{Total} & 163{,}2 \\\hline
\end{tabular}
\end{center}
\[ \ecm=\frac{163{,}2}{40}=\mathbf{4{,}08}. \]
\end{corbox}

\begin{corbox}{Corrigé 2 — Variance (formule directe)}
On élève les écarts au carré : $8{,}8^2=77{,}44$ ; $4{,}8^2=23{,}04$ ; $0{,}8^2=0{,}64$ ;
$3{,}2^2=10{,}24$ ; $7{,}2^2=51{,}84$. Pondérés :
\[
309{,}76+184{,}32+6{,}4+122{,}88+311{,}04=934{,}4,\qquad V=\frac{934{,}4}{40}=\mathbf{23{,}36}.
\]
\end{corbox}

\begin{corbox}{Corrigé 3 — Variance (König--Huygens)}
$\sum n_i x_i^{\,2}=4(4)+8(36)+10(100)+12(196)+6(324)=16+288+1000+2352+1944=5600$.
\[ V=\frac{5600}{40}-10{,}8^2=140-116{,}64=\mathbf{23{,}36}. \]
Même résultat que par la formule directe. \checkmark
\end{corbox}

\begin{corbox}{Corrigé 4 — Écart type et cohérence}
\[ \ect=\sqrt{23{,}36}=\mathbf{4{,}83}. \]
On a bien $\ect=4{,}83\geqslant\ecm=4{,}08$ : cohérent (le carré surpondère les grands écarts).
\end{corbox}

\begin{corbox}{Corrigé 5 — Même moyenne, dispersions différentes}
\textbf{Série $A$} ($m=10$) : écarts $-2,0,2$, carrés $4,0,4$ :
$V_A=\dfrac{8}{3}\approx2{,}67$, $\ect_A\approx1{,}63$.\\
\textbf{Série $B$} ($m=10$) : écarts $-8,0,8$, carrés $64,0,64$ :
$V_B=\dfrac{128}{3}\approx42{,}67$, $\ect_B\approx6{,}53$.\\
Même moyenne, mais $\ect_B\gg\ect_A$ : la série $\mathbf{B}$ est \textbf{beaucoup plus dispersée}.
\end{corbox}
\end{document}
